\chapter{相关工作}

本章综述与本文系统设计直接相关的五个方向：GRPO 算法、策略坍缩与防坍缩、自进化训练、奖励建模与奖励作弊、强化学习训练框架。

\section{GRPO 算法}

GRPO（Group Relative Policy Optimization）\cite{grpo} 对同一查询的一组 rollout 做组内奖励归一化来估计优势，省去了价值网络，是当前智能体强化学习的主力算法。近期工作\cite{grpo_dpo} 进一步从理论上揭示其方差缩减主要由每步 rollout 总数（查询数 $\times$ 组大小）驱动，为通过扩大采样规模提升训练信号质量提供了理论依据。本文采用标准 GRPO 作为策略优化算法，并通过超生-淘汰-选组机制扩大有效采样规模。

\section{策略坍缩与防坍缩}

多轮智能体强化学习暴露出策略多样性坍缩的 Echo Trap 现象\cite{echo_trap}：策略在自进化过程中逐渐丧失多样性，所有轨迹趋于一致，导致梯度信号消失、训练停滞。这一发现表明，自进化训练过程中的稳定性监控不可或缺。本文通过固定熵正则预防策略坍缩，并设计训练监控机制实时检测坍缩指标。

\section{自进化训练}

自进化是当前人工智能领域的重点研究方向。自我博弈微调（Self-Play Fine-Tuning）\cite{self_play} 让语言模型通过与自身的博弈持续自我提升，是自进化方向的代表性工作。强化微调（RFT）\cite{rft_forgetting} 的研究表明，基于强化学习的持续微调比监督微调更不易遗忘，为以强化学习驱动智能体持续进化提供了动机。本文的自进化系统在精神上与这些工作相通——基于历史经验自主产生新训练信号，但更进一步：通过跨步沙箱状态继承使智能体在上轮产出的物理基础上继续工作，而非仅从文本层面自我提升。

\section{奖励建模与奖励作弊}

以大模型为裁判（LLM-as-a-judge）\cite{instructgpt} 缓解了规则奖励覆盖面窄的问题，但若裁判看到的仍只是智能体自述的轨迹文本，奖励作弊空间依旧存在——模型会学会在轨迹中谎报完成，而裁判无从核实。如何让奖励锚定环境真实状态而非智能体叙事，是自进化训练中奖励可信化的核心问题。本文通过沙箱确定性状态差分使奖励锚定在环境客观变更上，从取证机制上抑制奖励作弊。

\section{强化学习训练框架}

verl（基于 HybridFlow）\cite{verl} 支持推理与训练的灵活编排，为大规模智能体强化学习提供了工程底座。但框架本身不提供数据自主产生、奖励可信取证、跨步状态继承等自进化所需的核心机制。本文以 verl 为底座，通过其官方注入点接入自定义采样与选组层，在框架之上构建自进化闭环。

\section{本章小结}

现有工作在算法层面已较成熟，但在"如何自主、可信地供给训练数据与奖励，并使系统持续自进化运转"这一系统与数据层面仍有明显空白。本文正是面向这一空白，以多智能体合成数据、差分驱动奖励、跨步自进化闭环，构建一套自进化强化学习系统。
